\section{Redes de puntuación condicionadas al ruido (NCSN)}

\subsection{Problema de estimación de puntuación en regiones de baja densidad}

Aprender directamente el score function de la distribución de datos original $p_{\text{data}}(x)$ mediante score matching tiene un problema grave.

\begin{warningbox}{Inexactitud de la puntuación en regiones de baja densidad}
En distribuciones multimodales, existen \textbf{regiones de baja probabilidad} (low density regions) entre los diferentes modos. En estas zonas:
\begin{enumerate}
    \item Hay muy pocos datos de entrenamiento, por lo que la estimación del score es extremadamente inexacta.
    \item Sin embargo, el punto de partida para el muestreo mediante Langevin dynamics suele caer en estas regiones de baja densidad (ya que la distribución inicial es usualmente uniforme o gaussiana).
    \item Una puntuación inexacta provoca que la fase inicial del proceso de muestreo esté llena de ruido; las partículas se mueven erráticamente ("huyen") en las zonas de baja densidad, resultando difícil llegar eficazmente a las regiones de alta densidad.
\end{enumerate}
\end{warningbox}

El ponente ilustra este problema con una animación de una distribución bidimensional: al partir de puntos de muestreo uniformes, si la puntuación en las zonas de baja densidad es inexacta, las partículas realizarán un movimiento aleatorio durante las etapas iniciales del muestreo en lugar de desplazarse de manera direccional hacia las regiones de alta densidad.

\subsection{Solución mediante ruido de múltiples escalas}

La idea central para resolver el problema de la inexactitud de la puntuación en zonas de baja densidad consiste en añadir \textbf{niveles distintos de ruido gaussiano} a los datos, con lo cual se "rellenan" las regiones de baja densidad.

Cuando la varianza del ruido $\sigma^2$ es grande, la distribución $q_\sigma(x)$ después de añadir el ruido se vuelve más ``suave'':
\begin{itemize}
    \item Las regiones de baja densidad originales quedan ``rellenas'' por el ruido; ya no existen zonas con probabilidad próxima a cero.
    \item El score function puede estimarse con precisión en todo el espacio.
    \item Pero el costo es que la distribución se vuelve borrosa, perdiendo información detallada.
\end{itemize}

Cuando la varianza del ruido $\sigma^2$ es pequeña:
\begin{itemize}
    \item La distribución después de añadir el ruido se acerca más a la distribución de datos originales e incluye información estructural fina.
    \item Pero la puntuación en las regiones de baja densidad sigue siendo inexacta.
\end{itemize}

\begin{importantbox}{La idea central de NCSN: aprendizaje de puntuaciones a múltiples escalas}
La idea central de NCSN (Noise-Conditioned Score Networks) es aprender simultáneamente el score function para varios niveles de ruido:
$$s_\theta(x, \sigma) \approx \nabla_x \log q_\sigma(x)$$
donde $\sigma$ toma una serie de valores decrecientes $\sigma_1 > \sigma_2 > \cdots > \sigma_L$. La red toma $x$ y $\sigma$ como entrada y devuelve la puntuación correspondiente al nivel de ruido.
\end{importantbox}

\subsection{Dinámica de Langevin con recocido (Annealed Langevin Dynamics)}

Con el score function a múltiples escalas, para el muestreo se utiliza \textbf{dinámica de Langevin con recocido} (Annealed Langevin Dynamics):

\begin{importantbox}{Algoritmo de muestreo mediante Dinámica de Langevin con recocido}
\begin{enumerate}
    \item Comenzar con gran varianza ($\sigma_1$ grande, distribución borrosa): la puntuación es bastante precisa en todo el espacio, guiando a las partículas desde su posición inicial hacia las regiones de alta densidad.
    \item Reducir progresivamente el nivel de ruido ($\sigma_2, \sigma_3, \ldots$): la puntuación se vuelve cada vez más precisa, refinando la guía de las partículas hacia las zonas de alta densidad de la distribución real.
    \item Converger en el nivel de ruido mínimo ($\sigma_L$ pequeño, cercano a la distribución original): las partículas se encuentran en las regiones de alta densidad de la distribución original.
\end{enumerate}
En cada nivel de ruido se ejecutan varios pasos de Langevin dynamics y luego se cambia al siguiente (nivel de ruido menor).
\end{importantbox}

\begin{knowledgebox}{Analogía física de la Dinámica de Langevin con recocido}
Este proceso es análogo al recocido (annealing) de metales:
\begin{itemize}
    \item Fase de alta temperatura (gran ruido): las moléculas pueden moverse libremente, cruzar barreras y explorar el panorama energético completo.
    \item Enfriamiento lento (reducción progresiva del ruido): las moléculas se estabilizan gradualmente en la configuración de menor energía.
    \item Fase de baja temperatura (poco ruido): las moléculas quedan bloqueadas en estados de baja energía precisos.
\end{itemize}
Análogo a la programación de temperaturas en el recocido simulado (simulated annealing).
\end{knowledgebox}

\subsection{Relación con el proceso inverso de los modelos de difusión}

El proceso de desruido a múltiples escalas mediante dinámica de Langevin con recocido guarda una relación profunda con el proceso inverso de los modelos de difusión:

\begin{importantbox}{Correspondencia entre Score-Based Model y DDPM}
\begin{center}
\begin{tabular}{lcc}
\toprule
\textbf{Concepto} & \textbf{Perspectiva DDPM} & \textbf{Perspectiva Score-Based} \\
\midrule
Proceso hacia adelante & Añadir ruido progresivamente (varianza creciente) & Distribución de datos se vuelve progresivamente borrosa \\
Proceso inverso & Desruido progresivo & Dinámica de Langevin con recocido \\
Salida de la red & Ruido $\epsilon_\theta$ & Puntuación $s_\theta$ \\
Objetivo de entrenamiento & Predicción de ruido L2 & Score matching desruidor \\
$\sigma_t$ grande & Inicios del proceso hacia adelante | Inicio del recocido (navegación gruesa) \\
$\sigma_t$ pequeño | Finales del proceso hacia adelante | Finales del recocido (refinamiento) \\
\bottomrule
\end{tabular}
\end{center}
El incremento progresivo de la varianza en el proceso hacia adelante de DDPM corresponde exactamente a la distribución a múltiples escalas desde poco ruido hasta mucho ruido en los modelos basados en puntuación; el proceso inverso de desruido de DDPM corresponde exactamente al refinamiento progresivo desde gran ruido hasta poco ruido en la dinámica de Langevin con recocido.
\end{importantbox}

\subsection{Resumen de la sección}

NCSN resuelve el problema de inexactitud en la estimación del score function en zonas de baja densidad mediante la introducción de ruido a múltiples escalas. La dinámica de Langevin con recocido comienza con gran ruido (navegación gruesa) y reduce progresivamente hacia poco ruido (refinamiento), lo que permite que el proceso de muestreo cruce eficazmente las zonas de baja densidad mientras converge finalmente a la distribución objetivo precisa. Este marco es matemáticamente totalmente equivalente al proceso inverso de desruido de DDPM, aunque ofrece una perspectiva e intuición distintas.
